\chapter{Planeamento do cinco inicial e distribuição de minutos}
\label{chap:rotacoes}

O módulo de planeamento de rotações de PLANTIA usa programação por restrições com o solver
OR-Tools CP-SAT. Conforme \path{docs/ROTATION_PLANNING.md}, seleciona
o cinco inicial e constrói uma rotação completa para os 48 minutos
regulamentares a partir dos modelos do módulo de avaliação de jogadores e alinhamentos.

\section{Entradas e função objetivo}

O plantel candidato é construído com o registo de equipa mais recente de
cada jogador nos 30 dias anteriores à data do cenário. As características
individuais usam apenas participações de dias anteriores, exigindo cinco
participações na temporada. A data deve ser posterior ao treino dos
modelos, e os hashes dos dados devem corresponder aos registados no
experimento de desempenho.

As posições provêm de uma fotografia biográfica sem versionamento
histórico; disponibilidade e pertença ao plantel são hipóteses a rever.
Candidatos sem posição ou histórico suficiente são excluídos. O adversário
não é explicitamente modelado; no modelo aditivo de titulares, a contribuição
de um adversário fixo seria constante entre alternativas.

Para previsões não negativas de pontos, ressaltos e assistências,
a prioridade individual por omissão é
\[
U_i=\widehat{P}_i+1{,}2\widehat{R}_i+1{,}5\widehat{A}_i.
\]
Previsões negativas são truncadas em zero. Sendo \(m_i\) os minutos,
\(s_i\) a indicação de titular e \(\beta_i\) o coeficiente aditivo do
módulo de avaliação de jogadores e alinhamentos, maximiza-se
\[
J=\sum_i U_i\frac{m_i}{48}+\lambda\sum_i\beta_i s_i,
\qquad \lambda=1.
\]
O intercepto de vantagem de casa é omitido por ser constante; jogadores
desconhecidos recebem coeficiente nulo e são sinalizados.

Os pesos exprimem preferências do planeamento. Como o estimador prevê
produção por participação, multiplicar pela fração de minutos constitui
uma heurística, não uma previsão de produção marginal por minuto.
A utilidade não mede pontos ou vitórias esperadas e pode favorecer
jogadores cuja produção histórica reflete maior utilização.

\section{Variáveis e restrições}

O modelo usa variáveis booleanas \(x_{irt}\), que indicam se o jogador
\(i\) ocupa o papel \(r\) no bloco \(t\), e \(y_{it}\), que indicam a sua
presença em campo. Com duração de bloco \(\Delta\),
\[
y_{it}=\sum_r x_{irt},\qquad m_i=\Delta\sum_t y_{it},
\qquad s_i=y_{i0}.
\]
Cada jogador ocupa no máximo um papel por bloco e só pode ocupar posições
para as quais é elegível. Por omissão, existem 12 blocos de quatro minutos,
cada um com dois lugares de base, dois de extremo e um de poste.
Esta composição é uma hipótese configurável, não uma regra obrigatória da NBA.

Em todos os blocos há exatamente cinco jogadores distintos. Se \(z_i\)
indica utilização na rotação, as restrições por omissão incluem
\[
8\leq\sum_i z_i\leq10,\qquad
8z_i\leq m_i\leq36z_i.
\]
A cobertura total implica \(\sum_i m_i=240\). Nenhum jogador pode disputar mais de três blocos consecutivos; cada janela
de quatro blocos satisfaz \(\sum_{t=k}^{k+3}y_{it}\leq3\).
Um bloco fora de campo permite regressar. As pausas entre períodos não
reiniciam esta contagem. Jogadores indisponíveis têm zero minutos e os
marcados como titulares obrigatórios devem integrar o primeiro bloco.

Os limites são configuráveis, com minutos múltiplos da duração dos blocos.
O modelo permite várias substituições na mesma fronteira e não representa
prolongamentos, faltas, lesões durante o jogo ou respostas do adversário.

\section{Resolução e verificação}

O comando \path{scripts/plan_rotation.py} exporta entradas editáveis,
solução, sequência de papéis por bloco, minutos e figura da rotação.
O solver utiliza coeficientes escalados por 1\,000 e arredondados, uma
semente de 42 e um único trabalhador para favorecer a reprodutibilidade.

O estado \texttt{OPTIMAL} demonstra optimalidade apenas para o objetivo
inteiro e as restrições fornecidas. \texttt{FEASIBLE} indica uma solução
válida sem prova de optimalidade no tempo disponível; \texttt{INFEASIBLE}
e \texttt{UNKNOWN} não produzem uma rotação bem-sucedida.
Uma verificação independente confirma posições, disponibilidade, cinco
inicial, minutos e descanso no horário exportado.

As entradas guardadas podem ser editadas para alterar limites e
indisponibilidades. Alterar apenas a data num ficheiro de configuração não
atualiza as previsões; um novo cenário temporal exige reconstruir as entradas.
